% ============================================================================
% Fichier  : partie2/P02_C06_preparer_planifier.tex
% Titre    : Préparer et Planifier la Mission
% Auteur   : MINKA MI NGUIDJOÏ Thierry Emmanuel
% Version  : 1.0 -- octobre 2026
% ============================================================================
% Sources : cours Audit_SI de l'auteur (travaux préparatoires, documents à
% solliciter, strates documentaires, lecture du SI en couches, entretiens par
% niveau décisionnel), mis en cohérence avec les articles A0 et A1 (cotation
% NEK, mesures µ, U, H) ; ordres de grandeur d'échantillonnage repris du
% cours GCI (chapitre 11), indicatifs et non normatifs.
% ============================================================================

\chapter{Préparer et Planifier la Mission}
\label{chap:planifier}

\epigraph{%
    « Le temps passé à préparer\\
    se reprend sur le terrain, au double. »%
}{Adage des chefs de mission}

\niveauChapitre{Fondamental}{%
    Maîtrise des Chapitres~\ref{chap:risque} (cartographie) et~\ref{chap:mission}
    (cycle, outils O0 à O9)%
}

\begin{boxobjectifs}
\begin{itemize}
    \item Distinguer les travaux préparatoires de la planification.
    \item Rédiger les termes de référence d'une mission en répondant à quatre
          questions~: quand, comment, par qui, avec quoi.
    \item Établir la liste des documents sollicités, et coter \NEK{} ou
          \NON{} un document absent selon le cas.
    \item Prendre connaissance d'un système d'information par couches, par
          strates documentaires et par niveaux d'interlocuteurs.
    \item Évaluer a priori le contrôle interne et produire la cartographie
          préliminaire des risques.
    \item Dimensionner les échantillons et rédiger le plan d'audit.
\end{itemize}
\end{boxobjectifs}

\bigskip

% ============================================================================
\section*{Question directrice : que faut-il savoir avant d'auditer~?}
% ============================================================================

AUDIX CONSEIL dispose de huit semaines pour auditer LiZbiZ~: quatre
centres de données, douze agences, une application intégrée, des milliers
de comptes et de mouvements de stock. Impossible de tout voir. M.~BEYALA
consacre donc les deux premières semaines à décider \emph{quoi} auditer,
\emph{où}, \emph{avec qui} et \emph{comment}, avant le moindre test.

Ce chapitre décrit ce travail. Il sépare la préparation de la
planification (section~\ref{sec:c06-deux-phases}), présente les termes de
référence (section~\ref{sec:c06-tdr}) et les documents sollicités
(section~\ref{sec:c06-documents}), puis la prise de connaissance
(section~\ref{sec:c06-connaissance}), l'évaluation a priori et la
cartographie préliminaire (section~\ref{sec:c06-risques}), et enfin le plan
d'audit (section~\ref{sec:c06-plan}).

\separateur

% ============================================================================
\section{Préparer, puis Planifier}
\label{sec:c06-deux-phases}
% ============================================================================

\begin{boxdef}
\termecle{Travaux préparatoires}\index{travaux préparatoires} (phase~0)~:
diligences qui conduisent à la lettre de mission et à la constitution de
l'équipe~: analyse préalable des risques, termes de référence, ressources.

\medskip
\termecle{Planification}\index{planification} (phase~1)~: diligences qui
préparent le terrain à partir de la lettre de mission~: prise de
connaissance, cartographie préliminaire des risques, décomposition de
l'espace d'audit, plan d'audit.
\end{boxdef}

On confond souvent les deux, parce que, du point de vue du contrat, la
mission commence à la signature de la lettre de mission. Du point de vue de
la démarche, la préparation \emph{produit} la lettre de mission~; la
planification l'\emph{exploite}. Les acteurs diffèrent aussi~: la
préparation associe étroitement le commanditaire, la planification
mobilise surtout l'équipe d'audit.

% ============================================================================
\section{Les Termes de Référence}
\label{sec:c06-tdr}
% ============================================================================

\subsection{Quatre questions}

Les termes de référence (outil~O2) répondent à quatre questions, qui
déterminent le coût et la faisabilité de la mission.

\begin{tableaucours}{L{2.4cm}L{5.0cm}L{4.8cm}}{%
    \tableauHeader{Question} &
    \tableauHeader{Contenu} &
    \tableauHeader{Chez LiZbiZ}
}
Quand~? & Période auditée, durée de la mission & Janvier à avril~2026~;
huit semaines \\
Comment~? & Type d'audit, objectifs, périmètre, référentiels, points
d'attention du commanditaire & Audit d'application (stocks) et audit de
sécurité (contrôle d'accès)~; COSO, ISO/IEC~27002 \\
Par qui~? & Profils requis, auditeurs internes ou externes, absence de
conflit d'intérêts & Un chef de mission, un auditeur réseau et sécurité, un
auditeur données (SQL) \\
Avec quoi~? & Moyens matériels, accès, budget & Accès en lecture à la base,
outil d'analyse de données, déplacements à Garoua et Bafoussam \\
\end{tableaucours}
\vspace{-0.8em}
\begin{center}{\small\itshape Tableau~6.1. Les quatre questions des termes
de référence}\end{center}

\subsection{Les livrables de la préparation}

À l'issue de la phase~0, le dossier de la mission contient les termes de
référence, la fiche technique (O1), la lettre de mission signée, et selon
le cas le contrat, l'engagement de confidentialité, la déclaration d'absence
de conflit d'intérêts de chaque auditeur et les autorisations d'accès aux
sites.

% ============================================================================
\section{Les Documents Sollicités}
\label{sec:c06-documents}
% ============================================================================

\subsection{Trois strates de documents}

La liste des documents sollicités (outil~O3) se construit selon le type
d'audit. Les documents n'ont pas tous la même valeur~: on les classe en trois
strates.

\begin{tableaucours}{L{2.4cm}L{3.8cm}L{3.0cm}L{3.0cm}}{%
    \tableauHeader{Strate} &
    \tableauHeader{Exemples} &
    \tableauHeader{Ce qu'ils révèlent} &
    \tableauHeader{Si absents}
}
1. Stratégique & Schéma directeur, politique de sécurité, plan de
continuité, rapports d'audit précédents & L'engagement de la direction &
Gouvernance informelle, recommandations non suivies \\
2. Organisationnelle & Organigramme de la DSI, cartographie des
applications et des interfaces, matrice des droits & Qui fait quoi &
Responsabilités floues, cumuls de fonctions \\
3. Opérationnelle & Journaux d'incidents, de sauvegardes, de traitements~;
inventaires & Ce qui se fait réellement & Impossible de distinguer le
déclaré du réel \\
\end{tableaucours}
\vspace{-0.8em}
\begin{center}{\small\itshape Tableau~6.2. Les strates documentaires}\end{center}

La confrontation des strates est en soi une source de constats~: une
politique de sécurité (strate~1) qui exige une revue trimestrielle des
droits, sans aucune trace de revue dans les journaux (strate~3), révèle
l'écart entre l'affiché et le pratiqué.

\subsection{Un document absent~: \NEK{} ou \NON{}~?}

Chaque document est sollicité avec une priorité (haute ou moyenne), un
interlocuteur et une date de remise. Lorsqu'un document n'est pas remis,
l'auditeur doit qualifier cette absence avec soin.

\begin{tableaucours}{L{4.8cm}cL{5.6cm}}{%
    \tableauHeader{Situation} &
    \tableauHeader{Cote} &
    \tableauHeader{Justification}
}
Le document n'est pas remis malgré une relance, sans qu'on sache s'il
existe & \NEK{} & L'information manque~: on ne peut pas conclure \\
L'audité confirme par écrit que le document n'existe pas, alors que le
critère l'exige & \NON{} & L'absence est établie~: le critère n'est pas
satisfait \\
Le document est remis, mais périmé ou non approuvé & \NON{} & Le critère
(document à jour et approuvé) n'est pas satisfait \\
\end{tableaucours}
\vspace{-0.8em}
\begin{center}{\small\itshape Tableau~6.3. Coter un document
manquant}\end{center}

\begin{boxlecture}{Documents sollicités}
Sur 21 documents sollicités le 5~mai, au 22~mai et après une relance~: 12
sont remis et conformes (\OUI{}), 3 sont remis mais périmés (\NON{}), 6 ne
sont pas remis (\NEK{}), dont le plan de continuité et la politique de
sécurité.

\medskip
\textbf{Constatation.} $\mu \approx 0{,}57$~; $U \approx 0{,}29$~;
$H \approx 1{,}38$~bit.

\medskip
\textbf{Analyse} (confiance moyenne). Avant tout test, plus d'un quart de
l'information de base manque, et deux des documents manquants relèvent de
la strate stratégique. Le plan d'audit doit réserver du temps à les
obtenir, ou à établir qu'ils n'existent pas~: chaque \NEK{} levé avant le
terrain réduit l'incertitude de toute la mission.

\medskip
\hfill\lecture{0{,}57}{0{,}29}{1{,}38}
\end{boxlecture}

\begin{boxalert}
\textbf{Violation typique.} Commencer les tests sans avoir obtenu la matrice
des droits, puis conclure sur la séparation des tâches à partir des seuls
entretiens.

\smallskip
\textbf{Règle~:} un document de priorité haute non remis est relancé par
écrit, et son absence est consignée dans le suivi des documents (O9). On ne
conclut jamais sur un contrôle dont la preuve principale manque~: il reste
\NEK{}.
\end{boxalert}

% ============================================================================
\section{Prendre Connaissance du Système d'Information}
\label{sec:c06-connaissance}
% ============================================================================

\subsection{Une prise de connaissance bornée et hiérarchisée}

\begin{boxdef}
\termecle{Prise de connaissance}\index{prise de connaissance}~: ensemble des
diligences par lesquelles l'auditeur acquiert, avant le terrain, une
compréhension suffisante de l'organisation, de son système d'information,
de ses processus et de ses risques pour définir un programme de travail
pertinent. Elle produit le \emph{dossier de prise de connaissance}, première
pièce du dossier d'audit.
\end{boxdef}

Elle est \emph{bornée} par le périmètre de la lettre de mission~: on ne
documente que ce qui le concerne. Elle est \emph{hiérarchisée}~: on lit les
documents par strates. Elle produit un livrable actionnable~: la
cartographie préliminaire des risques. Dans le vocabulaire de
l'auditologie, elle réduit l'incertitude initiale de l'auditeur jusqu'au
niveau qui permet de formuler des hypothèses de risque~\cite{minka2026auditologie}.

\subsection{Lire le SI en couches}

L'auditeur ne peut pas juger une application sans comprendre ce sur quoi
elle repose. Il lit le système d'information en quatre couches, de bas en
haut~:
\begin{enumerate}
    \item \textbf{Infrastructure}~: matériels, réseaux, hébergement~; ici,
          les quatre centres de données et les liaisons avec les agences.
    \item \textbf{Bases de données et intergiciels}~: la base de
          l'application intégrée, les interfaces.
    \item \textbf{Applications}~: l'application intégrée, ses paramétrages et
          ses contrôles.
    \item \textbf{Données et accès}~: profils, droits, journaux, conservation.
\end{enumerate}

\subsection{Interroger par niveau}

Les documents décrivent ce qui est formalisé~; les entretiens révèlent ce
qui est pratiqué. Les entretiens de prise de connaissance ne recueillent
pas encore de preuves~: ils situent les personnes, les tensions et les
problèmes connus.

\begin{tableaucours}{L{3.0cm}L{4.6cm}L{4.6cm}}{%
    \tableauHeader{Niveau} &
    \tableauHeader{Ce qu'on cherche} &
    \tableauHeader{Signal d'alerte}
}
Direction générale & Les objectifs assignés au SI, la tolérance au risque,
l'instance de pilotage & Le dirigeant ignore le rapport d'audit précédent \\
DSI et responsables informatiques & Organisation, parc applicatif,
prestataires, problèmes connus non résolus & « Tout va bien, nous n'avons
aucun incident » \\
Métiers et utilisateurs & La réalité vécue, les contournements, les
incidents non remontés & Un fichier tableur qui double l'application \\
\end{tableaucours}
\vspace{-0.8em}
\begin{center}{\small\itshape Tableau~6.4. Les entretiens de prise de
connaissance}\end{center}

\begin{boxmetier}{Auditeur SI}
\textbf{Conduire la prise de connaissance de LiZbiZ}
\begin{enumerate}
    \item Lire d'abord les documents de strate~1, puis~2, puis~3, et noter
          chaque écart entre strates.
    \item Rencontrer le directeur général, puis le responsable informatique,
          puis deux chefs d'agence et un magasinier.
    \item Demander à chacun les « problèmes connus non résolus »~: ce sont
          souvent les risques les plus sérieux.
    \item Cartographier les couches du SI, du centre de données au compte
          utilisateur.
    \item Consigner chaque entretien dans un compte rendu daté.
\end{enumerate}
\end{boxmetier}

% ============================================================================
\section{Évaluer a Priori et Cartographier}
\label{sec:c06-risques}
% ============================================================================

\subsection{L'évaluation a priori du contrôle interne}

Avant les tests, l'auditeur forme une première opinion sur le contrôle
interne, à partir des documents et des entretiens~: les contrôles
existent-ils, sont-ils conçus pour traiter les risques
(Chapitre~\ref{chap:controle-interne})~? Cette évaluation ne conclut rien~;
elle oriente. Un contrôle jugé bien conçu sera testé en efficacité~; un
contrôle absent ou mal conçu n'a pas besoin d'être testé~: c'est déjà un
constat, et l'auditeur portera ses tests directement sur les données.

\subsection{La cartographie préliminaire des risques}

L'auditeur part de la cartographie des risques de l'organisation, si elle
existe (Chapitre~\ref{chap:risque}). Sinon, il en établit une à partir de
la prise de connaissance. Il y ajoute ce que l'auditologie invite à
mesurer~: ce que l'on sait de chaque risque. On audite en priorité ce qui
est \emph{le plus risqué} et \emph{le moins bien connu}.

\begin{boxscenario}{Opération LIZBIZ. Choisir les zones d'audit}
\textbf{Constatations.} Le registre de Mme~NKENG classe trois risques comme
intolérables (ajustement du stock, fournisseurs fictifs, comptes d'agents
partis). La liste des documents laisse inconnus le plan de continuité et la
politique de sécurité. Les entretiens signalent un fichier tableur de suivi
des stocks tenu à Garoua en parallèle de l'application.

\medskip
\textbf{Analyse} (confiance moyenne). Le plan d'audit retient trois zones~:
la gestion des stocks (risque élevé, connaissance moyenne), le contrôle
d'accès, domaine~D (risque élevé, connaissance faible), et l'agence de
Garoua (contournement de l'application). La continuité est reportée à une
mission ultérieure, faute de temps, et le rapport le dira.
\end{boxscenario}

% ============================================================================
\section{Dimensionner et Rédiger le Plan d'Audit}
\label{sec:c06-plan}
% ============================================================================

\subsection{Dimensionner les échantillons}

L'étendue des tests dépend de la fréquence du contrôle testé et du risque
(Chapitre~\ref{chap:risque}). Les ordres de grandeur suivants sont
couramment pratiqués~; ce ne sont pas des exigences normatives, et chaque
organisation d'audit fixe sa méthode. Le Chapitre~\ref{chap:preuve} traite
de l'échantillonnage statistique.

\begin{tableaucours}{L{5.0cm}L{4.0cm}}{%
    \tableauHeader{Fréquence du contrôle} &
    \tableauHeader{Taille d'échantillon}
}
Annuelle & 1 \\
Trimestrielle & 2 \\
Mensuelle & 2 à 5 \\
Hebdomadaire & 5 à 15 \\
Quotidienne & 20 à 40 \\
Plusieurs fois par jour & 25 à 60 \\
\end{tableaucours}
\vspace{-0.8em}
\begin{center}{\small\itshape Tableau~6.5. Ordres de grandeur
d'échantillonnage (indicatifs)}\end{center}

On retient le haut de la fourchette lorsque le risque est élevé ou le
contrôle critique, le bas lorsqu'il est faible. Lorsque le risque de
détection doit être très bas, comme pour les stocks de LiZbiZ
(Chapitre~\ref{chap:risque}), on complète l'échantillon par des tests sur
l'ensemble des données.

\subsection{Le plan d'audit}

Le plan d'audit (outil~O6), ou rapport d'orientation, clôt la
planification. Il est la feuille de route du terrain.

\begin{boxprocedure}
\textbf{Le contenu du plan d'audit.} Pour qu'il puisse être approuvé, puis
servir de référence lors de la revue de la mission, le plan d'audit
indique~:
\begin{itemize}
    \item le contexte, les objectifs et le périmètre, rappelés de la lettre
          de mission~;
    \item les référentiels et les critères retenus~;
    \item l'évaluation a priori du contrôle interne et la cartographie
          préliminaire des risques~;
    \item les zones d'audit retenues et celles écartées, avec la raison~;
    \item le programme de travail par zone~: domaines décomposés, tests,
          tailles d'échantillon, preuves attendues~;
    \item le calendrier, l'équipe et l'affectation des tâches~;
    \item la date et la signature de son approbation par le commanditaire.
\end{itemize}
\end{boxprocedure}

\begin{boiteRemarque}{Le plan d'audit peut-il changer en cours de mission~?}
Oui, et il le doit lorsque le terrain révèle un risque imprévu~: les
ajustements de fin de mois à Douala (Chapitre~\ref{chap:fraude}) n'étaient
pas dans le plan initial d'AUDIX CONSEIL. Mais toute modification est
écrite, motivée et approuvée par le commanditaire~; sinon, le périmètre
réellement audité ne correspond plus à celui que le rapport annonce.
\end{boiteRemarque}

\separateur

\begin{boxresume}
\textbf{Ce qu'il faut retenir de ce chapitre~:}
\begin{itemize}
    \item \textbf{La préparation produit la lettre de mission~; la
          planification l'exploite.}
    \item \textbf{Les termes de référence} répondent à quatre questions~:
          quand, comment, par qui, avec quoi.
    \item \textbf{Les documents se lisent par strates}~: stratégique,
          organisationnelle, opérationnelle~; leur confrontation révèle
          l'écart entre l'affiché et le pratiqué.
    \item \textbf{Un document absent} est coté \NEK{} si l'on ne sait pas
          s'il existe, \NON{} si son absence est établie alors que le
          critère l'exige.
    \item \textbf{La prise de connaissance} lit le SI en couches et interroge
          par niveaux~; elle réduit l'incertitude initiale.
    \item \textbf{Le plan d'audit} retient les zones les plus risquées et les
          moins connues, dimensionne les tests et se fait approuver.
\end{itemize}
\end{boxresume}

\bigskip

\begin{boxexo}[title={\ding{45}\quad Exercice~6.1. Classer les documents \niveaubase}]
Classez dans l'une des trois strates les documents suivants, et indiquez ce
que révélerait leur absence~: le schéma directeur, la matrice des droits, le
journal des sauvegardes, le rapport de la dernière mission d'audit,
l'inventaire des licences, la charte d'utilisation du SI.
\end{boxexo}

\begin{boxexo}[title={\ding{45}\quad Exercice~6.2. Coter les absences \niveaumoyen}]
Pour chacune des situations suivantes, indiquez la cote (\OUI{}, \NON{} ou
\NEK{}) du contrôle « un plan de continuité approuvé existe pour le centre
de Garoua », et justifiez.
\begin{enumerate}[(a)]
    \item Le plan est remis~; il date de 2019 et n'a jamais été approuvé.
    \item Le responsable informatique dit qu'il existe, mais ne le remet pas
          malgré deux relances.
    \item Le directeur général écrit qu'aucun plan n'a jamais été rédigé.
    \item Le plan est remis, approuvé en 2025, et couvre Garoua.
\end{enumerate}
\end{boxexo}

\begin{boxexo}[title={\ding{45}\quad Exercice~6.3. Le plan d'audit de l'ANPP \niveauexpert}]
La Chambre des comptes prépare l'audit des systèmes d'information de l'ANPP
(Chapitre~\ref{chap:audit-auditologie}).
\begin{enumerate}[(a)]
    \item Rédigez les termes de référence en répondant aux quatre questions.
    \item Établissez la liste des dix documents à solliciter en priorité,
          classés par strate.
    \item Proposez trois zones d'audit et justifiez-les par le risque et par
          l'état de la connaissance.
    \item Dimensionnez l'échantillon du test de la revue mensuelle des droits
          d'accès au système de paiement.
\end{enumerate}
\end{boxexo}

\bigskip

\begin{boxinfo}
\textbf{Transition.} Le plan d'audit a désigné les zones à auditer. Il faut
maintenant les découper en éléments que l'auditeur peut réellement
vérifier. Le Chapitre~\ref{chap:decomposition} présente le gabarit canonique
de décomposition de l'espace d'audit et la fiche d'assertion.
\end{boxinfo}
